% !TEX root = ../Thesis.tex

\chapter{Progettazione e implementazione}
\label{ch:proposal}

Il sistema sviluppato e' una Progressive Web App orientata all'uso su smartphone. La sua funzione principale e' acquisire il video dell'utente, stimare la posa tramite MediaPipe e trasformare i landmark in eventi interpretabili: ripetizione valida, ripetizione non valida o movimento ancora in corso. L'implementazione concreta del prototipo e' contenuta nel repository locale del progetto \cite{projectRepository2026}.

\section{Requisiti}

I requisiti funzionali individuati sono:

\begin{itemize}
  \item selezione dell'esercizio tra squat, deadlift e pressa militare;
  \item selezione del lato osservato dalla fotocamera;
  \item acquisizione del flusso video dalla camera frontale;
  \item visualizzazione dello scheletro stimato sopra il video;
  \item conteggio di ripetizioni valide e non valide;
  \item visualizzazione del motivo della no-rep;
  \item storico delle ultime ripetizioni;
  \item esportazione dei dati in CSV;
  \item funzionamento come PWA.
\end{itemize}

I requisiti non funzionali sono:

\begin{itemize}
  \item esecuzione locale del modello, senza dipendenza runtime da CDN;
  \item interfaccia mobile-first;
  \item separazione tra logica di interfaccia, acquisizione video e classificazione;
  \item testabilita' della logica decisionale senza fotocamera;
  \item dichiarazione esplicita dei limiti tecnici.
\end{itemize}

\section{Architettura}

L'architettura e' composta da tre livelli principali:

\begin{enumerate}
  \item \textbf{Interfaccia utente}: componente React principale, schermata di setup, contatori, video, canvas e storico sessione.
  \item \textbf{Acquisizione e inferenza}: hook \texttt{usePose}, responsabile di inizializzare MediaPipe, avviare la camera, processare i frame e disegnare lo scheletro.
  \item \textbf{Logica cinematica}: modulo \texttt{repLogic.js}, contenente calcolo angoli, smoothing, macchine a stati e regole rep/no-rep.
\end{enumerate}

\begin{figure}[htbp]
\myfloatalign
\fbox{\begin{minipage}{0.88\textwidth}
\centering
Fotocamera $\rightarrow$ MediaPipe Pose Landmarker $\rightarrow$ Landmark corporei $\rightarrow$ Calcolo angoli $\rightarrow$ State machine esercizio $\rightarrow$ Rep / No-rep / CSV
\end{minipage}}
\caption{Pipeline logica del sistema sviluppato.}
\label{fig:pipeline}
\end{figure}

La separazione tra hook e logica consente di mantenere il codice piu' manutenibile. MediaPipe e' trattato come sorgente di landmark, mentre la classificazione delle ripetizioni e' indipendente dal rendering. Questo permette di scrivere test automatici sulla state machine usando landmark sintetici.

\section{Integrazione di MediaPipe}

MediaPipe Pose Landmarker viene inizializzato tramite \texttt{FilesetResolver.forVisionTasks} e \texttt{PoseLandmarker.createFromOptions}. Il modello e i file WASM sono serviti localmente dalle cartelle \texttt{public/models} e \texttt{public/wasm}. Tale scelta segue la documentazione ufficiale, che prevede il caricamento del modello da un percorso relativo o assoluto \cite{mediapipePoseWeb2026}, e migliora l'indipendenza del prototipo da risorse remote.

L'inferenza e' eseguita in modalita' \texttt{VIDEO} tramite \texttt{detectForVideo}. Per evitare inferenze duplicate sullo stesso frame, il sistema memorizza l'ultimo \texttt{currentTime} del video e processa un frame solo se il valore cambia. Questa ottimizzazione e' coerente con gli esempi ufficiali di MediaPipe per il web \cite{mediapipePoseWeb2026}.

\section{Calcolo degli angoli}

La funzione \texttt{calculateAngle} riceve tre landmark e calcola l'angolo nel landmark centrale. Per aumentare la stabilita' del segnale viene applicato uno smoothing esponenziale:

\begin{lstlisting}
export function smoothAngle(prev, current) {
  if (prev === null) return current;
  return (current * SMOOTHING.alpha) + (prev * SMOOTHING.beta);
}
\end{lstlisting}

Nel prototipo sono usati due angoli principali per volta: un angolo primario, che guida la state machine, e un angolo secondario, utile per feedback o controlli aggiuntivi. Ad esempio, nello stacco l'angolo primario e' quello dell'anca e il secondario quello del ginocchio; nella pressa militare il primario e' il gomito e il secondario e' l'inclinazione del tronco.

\section{Macchine a stati}

Ogni esercizio e' modellato come macchina a stati finita. Questa scelta rende esplicite le fasi del movimento e riduce il rischio di contare piu' volte la stessa ripetizione.

\begin{table}[htbp]
\myfloatalign
\begin{tabular}{lll}
\toprule
Esercizio & Stati principali & Evento finale \\
\midrule
Squat & STANDING, DESCENDING, ASCENDING & VALID\_REP / NO\_REP \\
Deadlift & STANDING, SETUP, LIFTING, LOCKED, DROPPING & VALID\_REP / NO\_REP \\
Pressa militare & STANDING, DESCENDING, ASCENDING & VALID\_REP / NO\_REP \\
\bottomrule
\end{tabular}
\caption{Stati principali delle macchine a stati implementate.}
\label{tab:states}
\end{table}

La funzione \texttt{processFrame} agisce da dispatcher: riceve esercizio, stato precedente, landmark e lato della camera, poi invoca la funzione specifica. Ogni funzione restituisce il nuovo stato, l'eventuale evento, gli angoli calcolati e un indicatore visuale utile per il disegno dello scheletro.

\section{Regole per lo squat}

Lo squat viene analizzato osservando l'angolo del ginocchio e la relazione verticale tra anca e ginocchio. La discesa viene rilevata quando l'angolo del ginocchio diminuisce rispetto alla posizione eretta. La profondita' viene considerata sufficiente quando l'anca scende sotto il ginocchio e l'angolo del ginocchio supera la soglia inferiore configurata. Se l'utente torna in estensione senza aver raggiunto tale condizione, la ripetizione e' classificata come no-rep per mancato superamento del parallelo.

Questa regola e' una semplificazione della valutazione reale dello squat. Le regole ufficiali richiedono che la superficie superiore delle gambe all'articolazione dell'anca scenda sotto la parte superiore delle ginocchia \cite{ipfRules2024}. Nel prototipo questa condizione e' approssimata usando le coordinate normalizzate di anca e ginocchio stimate da MediaPipe.

\section{Regole per il deadlift}

Nel deadlift il sistema osserva angolo dell'anca, angolo del ginocchio e posizione verticale del polso. Il polso e' usato come proxy della mano e quindi, in modo approssimato, del bilanciere. La ripetizione viene considerata valida quando anca e ginocchio superano le soglie di estensione. Se durante la fase di tirata il polso scende oltre una soglia rispetto al punto piu' alto raggiunto, il sistema genera una no-rep per discesa del bilanciere durante la tirata.

La regola e' ispirata al criterio tecnico secondo cui il bilanciere non dovrebbe muoversi verso il basso prima del completamento dell'alzata \cite{ipfRules2024}. Il prototipo non rileva direttamente il bilanciere, quindi la decisione e' da intendere come indicazione cinematica, non come giudizio ufficiale.

\section{Regole per la pressa militare}

La pressa militare usa tre controlli: range del gomito, inclinazione del tronco e flessione del ginocchio. La fase di discesa inizia quando il gomito si flette; la ripetizione e' valida solo se viene raggiunto un range minimo e poi il gomito torna in estensione. Se il tronco supera una soglia di inclinazione, viene segnalata iperlordosi lombare; se il ginocchio si flette oltre la soglia rispetto alla posizione iniziale, viene segnalato uso delle gambe.

Questa implementazione distingue una pressa stretta da una spinta assistita. Anche in questo caso la scelta e' volutamente interpretabile: non viene addestrato un classificatore, ma vengono esplicitate condizioni geometriche controllabili.

\section{Gestione del tracking}

Ogni regola verifica la visibilita' dei landmark necessari. Se durante una ripetizione il tracking viene perso per piu' di un secondo, il sistema produce una no-rep per occlusione. Tale comportamento e' conservativo: in assenza di dati affidabili, e' preferibile non validare una ripetizione.

\section{Interfaccia e dati esportati}

La schermata iniziale guida l'utente nella scelta dell'esercizio, del lato camera e del setup. Durante l'analisi vengono mostrati contatori, angoli, video e scheletro sovrapposto. Lo storico mostra le ultime ripetizioni e consente l'esportazione CSV. Ogni riga del CSV contiene timestamp, ora, esercizio, lato, esito, angoli principali, stato finale ed errori.

\section{Progressive Web App}

La build usa \texttt{vite-plugin-pwa} per generare manifest e service worker. Poiche' il file WASM di MediaPipe supera il limite predefinito di Workbox, la configurazione aumenta \texttt{maximumFileSizeToCacheInBytes}. Questo permette di includere le risorse necessarie al modello nella build PWA \cite{vitePwaDocs2026}. La PWA non e' solo una scelta estetica: su mobile consente un accesso piu' naturale e riduce la distanza tra prototipo web e applicazione installabile.
